% !TEX root = ../../../main.tex
% Block card: detection.pan_matrix_profile -- "Pan matrix profile" (Part I, the blocks chapter).
\cardfield{What it does}
Asks ``at which time scales does anything in this span repeat?'': it computes the matrix profile
at a ladder of window lengths, from 30 seconds to an hour by default, draws them as one picture,
and says which lengths repeat more than noise of the same kind would.

\cardfield{When to use it}
Before choosing a window for \blockref{detection.matrix_profile}, \blockref{detection.fluss} or a
motif search, when you do not know the scale of the events. The picture shows where in time the
repeating happens; the verdict says whether it beats chance.

\cardfield{What goes in, what comes out}
In goes a span of one channel. Out comes a picture: one row per window length, longest at the top,
time along the bottom, bright where the window starting at that moment has a close look-alike of
the same length elsewhere. With it comes a verdict in words: ``clears from A to B s; strongest gap at
C s'' (the band of lengths that repeat beyond noise) or ``no length stands out''.

\cardfield{Settings that matter}
\emph{max\_window\_s} first: a length is only judged when the span holds at least 14 windows of
it, so a 4-hour span can judge lengths up to about 17 minutes and an 8-hour span up to about 34.
\emph{min\_window\_s} and \emph{n\_lengths} set the rest of the ladder. \emph{null\_draws} is how
many shuffled copies each length is compared with; 3 is the measured minimum, and 0 gives the
picture without a verdict.

\cardfield{Watch out for}
``Strongest gap at C s'' is \emph{not} the length of the event, and not a window to hand to another
block: read the whole band and the picture. ``No length stands out'' means nothing repeats more than
this kind of noise, not that nothing happens. Detrend first. Two occurrences must have the same
length to match, and only shape counts, not size. It is slow: one full matrix profile per length,
plus one per shuffled copy.
